% Pizarra de T3, página 54.
\begin{ejer}[De premio: el objeto final]\label{ejer:t3-objeto-final}
Muestra que $\Spec\mathbb Z$ es el objeto final en la categoría de esquemas: para todo esquema $X$ existe un único morfismo $X\to\Spec\mathbb Z$. \hfill [Vakil 7.3.J]
\end{ejer}
\begin{sol}
Ya probamos en el \cref{ejer:t3-morfismos-hacia-afin} que hay una biyección
\[
\Hom_{\Sch}(X,\Spec\mathbb Z)
\cong\Hom_{\Ring}(\mathbb Z,\mathcal O_X(X)),
\]
pero $\Hom_{\Ring}(\mathbb Z,\mathcal O_X(X))$ consta de exactamente un morfismo, ya que $\mathbb Z$ es inicial en $\Ring$. Se sigue que existe un único $X\to\Spec\mathbb Z$.
\end{sol}
